\chapter{Applications of Sequential Circuits: Registers \& Counters}
\label{chap:sequential-applications}

\section{Introduction}
Registers and counters constitute the computational core of digital systems. While registers store and spatially/temporally transform binary words, counters generate deterministic numerical sequences, divide clock frequencies, and orchestrate microprocessor instruction cycles.

---

\section{Shift Registers}
An $N$-bit register consists of $N$ cascaded flip-flops capable of storing an $N$-bit binary word. A **Shift Register** is capable of shifting binary data laterally (left or right) upon successive clock pulses.

\subsection{The Four Primary Shift Register Topologies}

\begin{figure}[htbp]
\centering
\begin{tikzpicture}[scale=0.85, transform shape]
    % 4 D-FFs for SISO
    \foreach \x in {0, 2.8, 5.6, 8.4} {
        \draw[fill=boxnavybg, draw=brandnavy, line width=1.2pt, rounded corners=3pt] (\x, 0) rectangle (\x+1.8, 2.2);
        \node at (\x+0.9, 1.1) [font=\footnotesize\bfseries\color{brandnavy}] {D-FF};
        \node at (\x+0.3, 1.7) [font=\tiny] {D};
        \node at (\x+1.5, 1.7) [font=\tiny] {Q};
        \node at (\x+0.3, 0.4) [font=\tiny] {$>$};
    }
    % Interconnections
    \draw[->, >=Stealth, thick] (-1.0, 1.7) node[left, font=\small\bfseries] {Serial Data In} -- (0, 1.7);
    \draw[->, >=Stealth, thick] (1.8, 1.7) -- (2.8, 1.7);
    \draw[->, >=Stealth, thick] (4.6, 1.7) -- (5.6, 1.7);
    \draw[->, >=Stealth, thick] (7.4, 1.7) -- (8.4, 1.7);
    \draw[->, >=Stealth, thick] (10.2, 1.7) -- ++(1.0,0) node[right, font=\small\bfseries] {Serial Data Out};

    % Common Clock
    \draw[->, >=Stealth, line width=1.2pt, brandpurple] (-0.8, -0.6) node[left] {$CLK$} -- (9.0, -0.6);
    \foreach \x in {0.3, 3.1, 5.9, 8.7} {
        \draw[line width=1.2pt, brandpurple] (\x, -0.6) -- (\x, 0.4);
    }
\end{tikzpicture}
\caption{4-bit Serial-In Serial-Out (SISO) Shift Register schematic using D Flip-Flops.}
\label{fig:siso-register}
\end{figure}

\subsubsection{1. Serial-In Serial-Out (SISO)}
\begin{itemize}
    \item \textbf{Data Entry:} Single bit per clock cycle via a single input line.
    \item \textbf{Data Retrieval:} Single bit per clock cycle via the last flip-flop's output.
    \item \textbf{Timing Overhead:} For an $N$-bit word, it requires **$N$ clock pulses to fully load** and an additional **$N$ clock pulses to shift out**.
\end{itemize}

\subsubsection{2. Serial-In Parallel-Out (SIPO)}
\begin{itemize}
    \item \textbf{Data Entry:} Serial bit stream into the first stage.
    \item \textbf{Data Retrieval:} All $N$ bits are accessed simultaneously across parallel output pins ($Q_0, Q_1, \dots, Q_{N-1}$).
    \item \textbf{Application:} Serial-to-Parallel digital communications (e.g., UART receiver, SPI slave shift registers).
\end{itemize}

\subsubsection{3. Parallel-In Serial-Out (PISO)}
\begin{itemize}
    \item \textbf{Data Entry:} Loaded simultaneously onto all flip-flops using combinational steering gates controlled by a **Shift/$\overline{\text{Load}}$** signal.
    \item \textbf{Data Retrieval:} Shifted out serially one bit per clock cycle.
    \item \textbf{Application:} Parallel-to-Serial conversion (e.g., converting microprocessor data bus bytes into serial USB/UART transmissions).
\end{itemize}

\subsubsection{4. Parallel-In Parallel-Out (PIPO)}
\begin{itemize}
    \item \textbf{Data Entry \& Retrieval:} All bits are loaded in parallel and read out in parallel.
    \item \textbf{Timing Overhead:} Requires only **1 clock pulse** to store and update the entire $N$-bit register word.
    \item \textbf{Application:} General-purpose CPU data buffer registers and memory latch stages.
\end{itemize}

---

\section{Binary Counters}
A counter is a sequential state machine that transitions through a predefined sequence of binary states upon receiving clock pulses.

\begin{definition}[Counter Modulus (MOD-N)]
The modulus of a counter is the total number of unique states it cycles through before repeating its sequence. An $n$-flip-flop counter has a maximum modulus of $\text{MOD-}2^n$.
\end{definition}

\subsection{Asynchronous (Ripple) Counters}
In an Asynchronous Counter, the external clock signal is applied **only to the first (LSB) flip-flop**. Each subsequent flip-flop is clocked by the output transition of the preceding stage.

\subsubsection{1. 3-Bit Asynchronous Ripple UP Counter}
\begin{itemize}
    \item \textbf{Construction:} 3 JK (or T) flip-flops with $J = K = 1$ (toggle mode).
    \item \textbf{Clocking Rule:} Negative-edge-triggered ($\downarrow$) flip-flops clocked from the **True Output ($Q$)** of the preceding stage count in **UP mode** ($0 \rightarrow 1 \rightarrow 2 \rightarrow 3 \rightarrow 4 \rightarrow 5 \rightarrow 6 \rightarrow 7 \rightarrow 0$).
\end{itemize}

\begin{figure}[htbp]
\centering
\begin{circuitikz}[scale=0.85, transform shape]
    % FF0
    \draw (0,0) node[flipflop JK, scale=0.8] (ff0) {$\text{FF}_0$};
    \node at (ff0.bpin 1) [font=\tiny] {J};
    \node at (ff0.bpin 3) [font=\tiny] {K};
    \node at (ff0.bpin 2) [font=\tiny] {$>$};
    
    % FF1
    \draw (4,0) node[flipflop JK, scale=0.8] (ff1) {$\text{FF}_1$};
    
    % FF2
    \draw (8,0) node[flipflop JK, scale=0.8] (ff2) {$\text{FF}_2$};
    
    % Logic 1 tied to all J, K
    \draw (-1.5, 1.2) node[above] {$V_{CC} (\mathbf{1})$} -- (0.8, 1.2);
    
    % External Clock
    \draw (-1.5, 0) node[left] {$CLK$} -- (ff0.bpin 2);
    
    % Ripple clock connections
    \draw (ff0.bpin 6) -- ++(0.5,0) -- ++(0,-1.0) -| (ff1.bpin 2) node[pos=0.2, above, font=\scriptsize] {$Q_0$};
    \draw (ff1.bpin 6) -- ++(0.5,0) -- ++(0,-1.0) -| (ff2.bpin 2) node[pos=0.2, above, font=\scriptsize] {$Q_1$};
    
    % Outputs
    \draw (ff0.bpin 6) -- ++(0.5,0) -- ++(0,1.2) node[above, font=\bfseries] {$Q_0\text{ (LSB)}$};
    \draw (ff1.bpin 6) -- ++(0.5,0) -- ++(0,1.2) node[above, font=\bfseries] {$Q_1$};
    \draw (ff2.bpin 6) -- ++(0.8,0) node[right, font=\bfseries] {$Q_2\text{ (MSB)}$};
\end{circuitikz}
\caption{Logic circuit of a 3-Bit Asynchronous (Ripple) UP Counter.}
\label{fig:async-up-counter}
\end{figure}

\subsubsection{2. Asynchronous DOWN Counter}
\begin{itemize}
    \item Clocking negative-edge-triggered flip-flops from the **Inverted Output ($\overline{Q}$)** of the preceding stage results in a **DOWN Counter** ($7 \rightarrow 6 \rightarrow 5 \rightarrow 4 \rightarrow 3 \rightarrow 2 \rightarrow 1 \rightarrow 0 \rightarrow 7$).
\end{itemize}

\subsubsection{3. Truncated Modulo-$N$ Asynchronous Counters}
To design a counter with modulus $N < 2^n$ (such as a **MOD-10 Decade / BCD Counter**):
\begin{enumerate}
    \item Determine required flip-flops: $2^n \ge N \implies 2^4 \ge 10 \implies n = 4$.
    \item Identify the binary state $N$ where the counter must reset: $10_{10} = 1010_2$ ($Q_3 = 1, Q_2 = 0, Q_1 = 1, Q_0 = 0$).
    \item Connect $Q_3$ and $Q_1$ to a **2-input NAND gate**, and feed its output to the active-low **$\overline{\text{CLR}}$ (Clear)** inputs of all four flip-flops.
    \item When count reaches $1010_2$, the NAND output goes LOW, immediately forcing all flip-flops back to $0000_2$.
\end{enumerate}

\begin{commonmistake}[Cumulative Propagation Delay in Ripple Counters]
In an $N$-bit ripple counter, each flip-flop introduces a propagation delay $t_{pd}$. The total time required for all bits to settle is $T_{total} = N \cdot t_{pd}$. To ensure correct counting without spurious output glitches:
\begin{equation}
    T_{clock} \ge N \cdot t_{pd} \implies f_{max} \le \frac{1}{N \cdot t_{pd}}
\end{equation}
For high-frequency counting, **Synchronous Counters** must be used instead.
\end{commonmistake}

---

\section{Synchronous Counters}
In a **Synchronous Counter**, the master clock pulse is applied simultaneously to the clock inputs of all flip-flops, eliminating cumulative ripple delay.

\subsection{3-Bit Synchronous UP Counter}
\begin{itemize}
    \item \textbf{Toggle Logic:} A stage toggles on clock pulse only if **all previous lower-order bits are $1$**.
    \begin{align*}
        T_0 &= 1 \quad (\text{FF}_0 \text{ toggles on every clock cycle}) \\
        T_1 &= Q_0 \quad (\text{FF}_1 \text{ toggles when } Q_0 = 1) \\
        T_2 &= Q_0 \cdot Q_1 \quad (\text{FF}_2 \text{ toggles when both } Q_0 = 1 \text{ and } Q_1 = 1)
    \end{align*}
\end{itemize}

\begin{figure}[htbp]
\centering
\begin{circuitikz}[scale=0.85, transform shape]
    % FFs
    \draw (0,0) node[flipflop T, scale=0.8] (t0) {$\text{FF}_0$};
    \draw (4,0) node[flipflop T, scale=0.8] (t1) {$\text{FF}_1$};
    \draw (8,0) node[flipflop T, scale=0.8] (t2) {$\text{FF}_2$};
    
    % T inputs
    \draw (t0.bpin 1) -- ++(-0.5,0) node[left] {$\mathbf{1}\text{ (5V)}$};
    \draw (t0.bpin 6) -- (t1.bpin 1) node[midway, above, font=\scriptsize] {$Q_0$};
    
    % AND gate for T2
    \node (and1) at (6.5, 0.8) [and gate US, draw, logic gate inputs=nn, scale=0.8] {};
    \draw (t0.bpin 6) ++(0.5,0) node[circ] {} |- (and1.input 1);
    \draw (t1.bpin 6) ++(0.5,0) node[circ] {} |- (and1.input 2);
    \draw (and1.output) -- (t2.bpin 1);
    
    % Common Clock Line
    \draw[line width=1.2pt, brandpurple] (-1,-1.5) node[left] {$CLK$} -- (8.3,-1.5);
    \draw[->, >=Stealth, line width=1.2pt, brandpurple] (0.3,-1.5) -- (t0.bpin 2);
    \draw[->, >=Stealth, line width=1.2pt, brandpurple] (4.3,-1.5) -- (t1.bpin 2);
    \draw[->, >=Stealth, line width=1.2pt, brandpurple] (8.3,-1.5) -- (t2.bpin 2);
\end{circuitikz}
\caption{Complete circuit schematic of a 3-Bit Synchronous Parallel UP Counter.}
\label{fig:sync-up-counter}
\end{figure}

\begin{example}[Design of a Synchronous Mod-6 Counter using JK Flip-Flops]
Design a Synchronous Counter that counts through the sequence: $0 \rightarrow 1 \rightarrow 2 \rightarrow 3 \rightarrow 4 \rightarrow 5 \rightarrow 0$ using JK Flip-Flops.

\tcblower
\textbf{Step 1: Determine Number of Flip-Flops}
Number of states $N = 6 \implies 2^n \ge 6 \implies n = 3$ Flip-Flops ($Q_2, Q_1, Q_0$).
Unused states: $6 (110_2)$ and $7 (111_2)$.

\textbf{Step 2: State Table and JK Excitation Mapping}
\begin{center}
\small
\begin{tabular}{ccc|ccc|cc|cc|cc}
\toprule
\multicolumn{3}{c|}{\textbf{Present State}} & \multicolumn{3}{c|}{\textbf{Next State}} & \multicolumn{2}{c|}{\textbf{$\text{FF}_2$}} & \multicolumn{2}{c|}{\textbf{$\text{FF}_1$}} & \multicolumn{2}{c}{\textbf{$\text{FF}_0$}} \\
$Q_2$ & $Q_1$ & $Q_0$ & $Q_2^+$ & $Q_1^+$ & $Q_0^+$ & $J_2$ & $K_2$ & $J_1$ & $K_1$ & $J_0$ & $K_0$ \\
\midrule
0 & 0 & 0 & 0 & 0 & 1 & 0 & X & 0 & X & 1 & X \\
0 & 0 & 1 & 0 & 1 & 0 & 0 & X & 1 & X & X & 1 \\
0 & 1 & 0 & 0 & 1 & 1 & 0 & X & X & 0 & 1 & X \\
0 & 1 & 1 & 1 & 0 & 0 & 1 & X & X & 1 & X & 1 \\
1 & 0 & 0 & 1 & 0 & 1 & X & 0 & 0 & X & 1 & X \\
1 & 0 & 1 & 0 & 0 & 0 & X & 1 & 0 & X & X & 1 \\
\bottomrule
\end{tabular}
\end{center}

\textbf{Step 3: K-Map Simplifications (with Don't Cares at 6 and 7):}
\begin{itemize}
    \item For $\text{FF}_2$: $\mathbf{J_2 = Q_1 Q_0}, \quad \mathbf{K_2 = Q_0}$
    \item For $\text{FF}_1$: $\mathbf{J_1 = \overline{Q}_2 Q_0}, \quad \mathbf{K_1 = Q_0}$
    \item For $\text{FF}_0$: $\mathbf{J_0 = 1}, \quad \mathbf{K_0 = 1}$
\end{itemize}

\textbf{Step 4: Lock-Out / Self-Starting Verification}
\begin{itemize}
    \item \textbf{From Unused State $6 (110_2)$:}
    \begin{align*}
        J_2=0, K_2=0 &\implies Q_2^+ = 1 \\
        J_1=0, K_1=0 &\implies Q_1^+ = 1 \\
        J_0=1, K_0=1 &\implies Q_0^+ = 1 \implies \text{Next State is } \mathbf{7 \, (111_2)}
    \end{align*}
    \item \textbf{From Unused State $7 (111_2)$:}
    \begin{align*}
        J_2=1, K_2=1 &\implies Q_2^+ = 0 \\
        J_1=0, K_1=1 &\implies Q_1^+ = 0 \\
        J_0=1, K_0=1 &\implies Q_0^+ = 0 \implies \text{Next State is } \mathbf{0 \, (000_2)}
    \end{align*}
\end{itemize}
The counter automatically recovers to the valid counting cycle within 2 clock pulses $\implies$ **Self-Starting Counter**.
\end{example}

---

\section{Shift Register Counters (Ring \& Johnson Counters)}

\subsection{Ring Counter}
A Ring Counter is a circular shift register where the output of the final stage is fed directly to the input of the first stage ($Q_{\text{last}} \rightarrow D_{\text{first}}$).
\begin{itemize}
    \item **Modulus:** $\text{MOD-}N$ ($N$ total states for $N$ flip-flops).
    \item **Initialization:** Must be preset with a single $\mathbf{1}$ (e.g., $1000_2$).
    \item **Counting Cycle (4-bit):**
    \begin{equation}
        1000 \longrightarrow 0100 \longrightarrow 0010 \longrightarrow 0001 \longrightarrow 1000 \longrightarrow \cdots
    \end{equation}
    \item **Decoding:** Inherent (each flip-flop output directly indicates a specific sequence state without external decoding gates).
\end{itemize}

\subsection{Johnson (Twisted Ring) Counter}
A Johnson Counter connects the **inverted output** of the final stage to the input of the first stage ($\overline{Q}_{\text{last}} \rightarrow D_{\text{first}}$).
\begin{itemize}
    \item **Modulus:** $\text{MOD-}2N$ ($2N$ total states for $N$ flip-flops).
    \item **Initialization:** Typically cleared to all zeros ($0000_2$).
    \item **Counting Cycle (4-bit, 8 States):**
    \begin{equation}
        0000 \rightarrow 1000 \rightarrow 1100 \rightarrow 1110 \rightarrow 1111 \rightarrow 0111 \rightarrow 0011 \rightarrow 0001 \rightarrow 0000
    \end{equation}
    \item **Decoding:** Decoded using simple 2-input AND gates ($S_0 = \bar{Q}_3\bar{Q}_0, S_1 = Q_0\bar{Q}_1, \dots$).
\end{itemize}

\begin{table}[htbp]
\centering
\small
\begin{tabularx}{\textwidth}{lXXX}
\toprule
\textbf{Counter Feature} & \textbf{Standard Binary Counter} & \textbf{Ring Counter} & \textbf{Johnson Counter} \\
\midrule
\textbf{Number of States} & $2^N$ & $N$ & $2N$ \\
\textbf{Clock Synchronization} & Asynchronous or Synchronous & Synchronous & Synchronous \\
\textbf{Decoding Logic} & $N$-input AND gates & None (Direct outputs) & 2-input AND gates \\
\textbf{Unused States (4-bit)} & 0 & 12 unused states & 8 unused states \\
\textbf{Primary Applications} & Arithmetic counting, timebases & Multiphase clocking, stepper motors & Frequency synthesis, timing generators \\
\bottomrule
\end{tabularx}
\caption{Comprehensive Comparison of Counter Architectures.}
\label{tab:counter-comparison}
\end{table}

---

\section{Summary \& Chapter Review}
\begin{quickrevision}[Unit 6 Takeaways]
\begin{itemize}
    \item \textbf{Shift Register Topologies:} SISO ($N$ in, $N$ out), SIPO (serial to parallel), PISO (parallel to serial), PIPO (instantaneous 1-clock buffer).
    \item \textbf{Ripple Counters:} Asynchronous clocking; simple hardware but limited by cumulative propagation delay ($f_{max} \le \frac{1}{N \cdot t_{pd}}$).
    \item \textbf{Synchronous Counters:} Clock connected to all flip-flops simultaneously; high operating speed.
    \item \textbf{Ring Counter:} $N$ flip-flops $\implies N$ states; circular shift of a single bit $\mathbf{1}$.
    \item \textbf{Johnson Counter:} $N$ flip-flops $\implies 2N$ states; inverted feedback from $\bar{Q}_{last}$ to $D_{first}$.
\end{itemize}
\end{quickrevision}
